\documentclass[12pt]{article}
\usepackage[margin=1in]{geometry}
\usepackage[colorlinks=true,citecolor=blue,linkcolor=blue,urlcolor=blue]{hyperref}
\usepackage{xcolor}
\usepackage{enumitem}
\linespread{1.3}

\newcommand{\referee}[1]{\medskip\noindent\textbf{Comment:} \textit{#1}\medskip}
\newcommand{\response}[1]{\noindent\textbf{Response:} #1}

\title{Response to Referee Report (R1)\\[6pt]
\large International Journal of Industrial Organization\\[4pt]
\large ``Frequent Losers as Cover Bidders in Public Procurement''\\
\large Genicolo-Martins \& Furquim de Azevedo}
\author{}
\date{March 2026}

\begin{document}
\maketitle

\noindent Dear Editor and Referee,

Thank you for the constructive and detailed report. We appreciate the
careful reading and the recommendation of Minor Revision. All four major
items (M1--M4) and all eight minor items have been addressed through
textual revisions. No new estimation was required; all changes are
editorial, consistent with the Minor Revision designation. Below we
describe each change, quoting the revised text so the referee can verify
without re-reading the full manuscript.

\bigskip

% ==========================================================================
\section*{Major Comments}
% ==========================================================================

\subsection*{M1: Prediction 5 contradicts Bajari--Ye tender-FE result}

\referee{Prediction 5 predicts FL residuals show ``magnitude larger than
non-FL residuals,'' but Section~7.4 reports the opposite under tender FE
(FL 0.38 < non-FL 0.86). The prediction contradicts the finding.}

\response{We revised Prediction~5 to target the exchangeability violation
without making a directional claim about residual magnitudes. The revised
prediction reads:}

\begin{quote}
\textbf{Prediction 5 (Bid Heterogeneity):} Under cover bidding, FL
firms' bid distributions are non-exchangeable with those of non-FL firms
conditional on observables. Specifically, the estimated first-stage
residuals should differ systematically across FL and non-FL firms,
consistent with firms playing distinct equilibrium strategies rather than
drawing from a common distribution. Cover bids---whether complementary or
coordinated---share a common signal (the designated winning bid) that
induces positive residual covariance among FL bidders, though its
magnitude relative to non-FL residual covariance depends on the
information environment.
\end{quote}

\noindent We also updated the cross-reference in Section~7.4 (Bajari--Ye
results) to be consistent with the revised prediction, replacing the
``stronger correlation'' language with ``systematically different
correlation \ldots consistent with distinct bidding processes.''

\noindent\textit{Location: Section~3.3 (sec\_conceptual.tex) and
Section~7.4 (sec\_results.tex).}

% --------------------------------------------------------------------------
\subsection*{M2: Balance test discussion inconsistency}

\referee{Section~6.2 (threats) honestly acknowledges that the pattern of
significant imbalance across all four observables is concerning. Section~7.2
contradicts this by dismissing imbalance as ``well inside negligible.''}

\response{We harmonized the two passages. In Section~7.2 (IV robustness),
we replaced the dismissive language with:}

\begin{quote}
The balance tests reveal statistically significant differences across all
four observable characteristics, a pattern we discuss as a potential
threat in Section~6.2. While the economic magnitudes are modest (all
standardized differences below $0.03\sigma$), we do not dismiss this
imbalance; rather, we rely on the sensitivity analysis and matching
estimates to bound the extent to which unobserved selection could explain
the price effect.
\end{quote}

\noindent In Section~6.2 (threats to exclusion restriction), we
strengthened the honest assessment:

\begin{quote}
Table~B.8 regresses pre-determined observables on the instrument. The
pattern of statistically significant imbalance across all four
observables is more concerning than any single coefficient; we rely on
the sensitivity analysis (Section~9) and matching estimates to assess
whether this imbalance is sufficient to explain the price effect.
\end{quote}

\noindent\textit{Location: Sections~6.2 (sec\_empirical\_strategy.tex)
and~7.2 (sec\_results.tex).}

% --------------------------------------------------------------------------
\subsection*{M3: Enriched Bajari--Ye first stage and Section~6.4 correction}

\referee{The enriched first stage (firm age + CNAE sector, $R^2$ increases
from 0.776 to 0.803, KS statistic unchanged) exists as estimation output
but is not described in the text. Also, Section~6.4 incorrectly lists
``reference price'' among first-stage controls.}

\response{\textbf{Part A.} We added a paragraph to Section~7.4
(Bajari--Ye results):}

\begin{quote}
\textbf{Enriched first stage.} To assess whether the FL--non-FL residual
differences reflect observable firm heterogeneity rather than strategic
behavior, we re-estimate the first-stage bid equation adding firm age and
CNAE sector dummies as controls. The $R^2$ increases modestly from 0.776
to 0.803, confirming that the original specification captures most
systematic variation in bid levels. Crucially, the
Kolmogorov--Smirnov statistic---the key test outcome for
exchangeability---is unchanged, ruling out that the FL--non-FL bid
distribution differences are driven by omitted firm-level controls.
\end{quote}

\noindent\textbf{Part B.} In Section~6.4 (Bajari--Ye test design), we
corrected the first-stage description, removing ``reference price'' and
explaining that it is absorbed by item FE:

\begin{quote}
The first stage regresses $\log$ bid values on observable cost shifters
(firm size, firm age, CNAE sector dummies, number of bids, and item and
year fixed effects) to obtain bid residuals. Reference price is excluded
as it varies at the tender level and is fully absorbed by the item fixed
effects.
\end{quote}

\noindent\textit{Location: Sections~6.4 (sec\_empirical\_strategy.tex)
and~7.4 (sec\_results.tex).}

% --------------------------------------------------------------------------
\subsection*{M4: Figure D.2 caption}

\referee{Figure~D.2 caption reads ``OLS (lower bound) and IV (upper
bound)'' which contradicts the text's framing of the welfare calculation
as ``back-of-the-envelope.''}

\response{We replaced the caption with:}

\begin{quote}
Welfare loss bounds computed using OLS and IV price-effect estimates as
back-of-the-envelope benchmarks. These bounds do not constitute causal
welfare estimates; see Section~9.5 for the full discussion of
assumptions. OECD benchmark range for cartel overcharges shown as dashed
lines.
\end{quote}

\noindent\textit{Location: Appendix~D (sec\_appendix.tex), Figure~D.2.}

\bigskip

% ==========================================================================
\section*{Minor Comments}
% ==========================================================================

\subsection*{minor-1: Remove version-tracking LaTeX tags}

\referee{Drafting comments (CO-AUTHOR EDIT tags) should be removed from
the final manuscript.}

\response{All 24 occurrences of \texttt{\% CO-AUTHOR EDIT: ...} comments
have been removed from 7 \texttt{.tex} files. No structural code was
affected---only inline drafting annotations were stripped.}

% --------------------------------------------------------------------------
\subsection*{minor-2: Section~6.4 first-stage variable description}

\referee{The first-stage description should be consistent with the
enriched specification.}

\response{Addressed as part of M3b above. Sections~6.4 and~7.4 now
describe the same set of first-stage controls.}

% --------------------------------------------------------------------------
\subsection*{minor-3: Learning / entry timing caveat}

\referee{The one-year age difference between FL and non-FL firms could
reflect entry timing across product markets rather than differential
learning.}

\response{We added a new subsection (``Learning and Entry Timing'') to
Section~8 (Mechanisms):}

\begin{quote}
Descriptive statistics show that FL firms are on average approximately
one year younger than non-FL firms at the time of their first BEC
participation. This age gap is consistent with the cover bidding
hypothesis (recently formed shell firms), but could also reflect
differential learning: younger firms may not yet have acquired the
market-specific knowledge needed to win competitively. We note that the
one-year age gap may also reflect firm entry timing across product
markets rather than learning per se; disentangling these channels would
require firm-level data on market entry sequencing, which is outside the
scope of this paper.
\end{quote}

\noindent\textit{Location: Section~8 (sec\_mechanisms.tex), new
subsection before ``Interpretation of the Convite--Preg\~ao
Difference.''}

% --------------------------------------------------------------------------
\subsection*{minor-4: Table cutoff rationale (HHI median, repeat partners)}

\referee{The HHI median split and repeat-partners $\geq 2$ cutoffs in the
network classification lack explicit economic rationale.}

\response{We expanded the footnote accompanying the network-split results
(Section~7.3) to include:}

\begin{quote}
Markets above the HHI sample median are classified as concentrated; the
median is used to avoid ad hoc threshold selection while maintaining
sample balance. Firms appearing together in at least two distinct tenders
are classified as having established repeat relationships; this threshold
minimizes classification error for rare pairings while capturing
sustained co-participation.
\end{quote}

\noindent\textit{Location: Section~7.3 (sec\_results.tex), footnote to
network-split table discussion.}

% --------------------------------------------------------------------------
\subsection*{minor-5: Imhof proxy coarseness caveat}

\referee{The comparison with Imhof-style screens uses a coarser
approximation than the full Imhof (2017) implementation.}

\response{We added a footnote to the literature review's bid-rigging
screens paragraph:}

\begin{quote}
Our comparison with Imhof-style screens (Section~9) uses a CV median
split as a coarser approximation of the full Imhof et al.\ (2019)
implementation, which employs multiple bid-level features and a trained
classifier. Our horse-race therefore provides a conservative lower bound
on the incremental value of the FL screen relative to a fully implemented
Imhof classifier.
\end{quote}

\noindent\textit{Location: Section~2 (sec\_literature.tex), footnote in
bid-rigging screens paragraph.}

% --------------------------------------------------------------------------
\subsection*{minor-6: IQR footnote economic framing}

\referee{The IQR threshold footnote should include an economic
interpretation, not just a statistical description.}

\response{We augmented the existing footnote in Section~4.2:}

\begin{quote}
Economically, this cutoff identifies firms with participation frequency
so extreme as to be difficult to rationalize as competitive optimization:
at this threshold, the marginal cost of bid preparation would exceed any
plausible expected return from participation under standard assumptions
about win probabilities.
\end{quote}

\noindent The original statistical description is preserved alongside
the new economic framing.

\noindent\textit{Location: Section~4.2 (sec\_data.tex), IQR footnote.}

% --------------------------------------------------------------------------
\subsection*{minor-7: FE-residualized gap in Introduction}

\referee{The FE-residualized gap (0.050 log points) is an important
finding that strengthens H1 and should appear in the Introduction.}

\response{We added to the first contribution paragraph:}

\begin{quote}
After absorbing item and time fixed effects, the FL price gap narrows to
0.050 log points (5.0\%), confirming that the unconditional estimate is
not driven by systematic differences in the types of goods or services
procured in FL-present tenders.
\end{quote}

\noindent\textit{Location: Section~1, first contribution paragraph
(sec\_introduction.tex).}

% --------------------------------------------------------------------------
\subsection*{minor-8: Compress LATE interpretation paragraph}

\referee{The LATE interpretation paragraph is standard and takes more
space than warranted given the IV is now supplementary evidence.}

\response{We compressed the paragraph from 6 lines to 3 sentences:}

\begin{quote}
The IV estimate of 21\% ($F = 396$) identifies a LATE for complier
tenders; exclusion restriction concerns are discussed in Section~6.2. We
treat the IV as supplementary evidence that brackets the OLS estimate
from above.
\end{quote}

\noindent\textit{Location: Section~7.2 (sec\_results.tex), LATE
interpretation paragraph.}

\bigskip
\bigskip

% ==========================================================================
\section*{Closing}
% ==========================================================================

We believe all referee comments have been addressed. The manuscript
compiles cleanly with no undefined references or warnings. We are
grateful for the quality of the review, which has improved the internal
consistency of the paper. The manuscript is ready for final decision.

\bigskip
\noindent Sincerely,\\
Darcio Genicolo-Martins \& Paulo Furquim de Azevedo\\
INSPER, S\~ao Paulo

\end{document}
